\section{Introduction}\label{sec:intro}

A central empirical puzzle in contemporary AI safety is \emph{emergent misalignment} (EM):
finetuning a large language model on a narrow, seemingly harmless dataset---insecure code,
or bad health advice---can cause the model to become broadly misaligned across unrelated
domains \cite{betley2025em,wang2025persona}. The phenomenon is fragile and regime-dependent.
The same intervention sometimes produces a null result (robust alignment), sometimes a gradual
drift, and sometimes a sudden behavioral collapse \cite{soligo2026em,nghiem2026trait}. There is,
at present, no predictive theory for \emph{when} an alignment intervention catastrophically
fails, nor for the \emph{order} of the failure (continuous versus abrupt), nor for early-warning
signals that precede it.

This is precisely the kind of question that the mathematics of \emph{critical transitions}
was built to answer. Bifurcation theory of fast--slow systems \cite{kuehn2011}, dynamic
transition theory for gradient flows \cite{mawang2008}, and the variational Free Energy
Principle (FEP) of theoretical biology \cite{friston2021} together supply order parameters,
threshold criteria, normal forms, and divergence laws. Yet this machinery has never been
applied to alignment, and the empirical EM literature uses none of its language. The two
bodies of work have remained disjoint. The contribution of this paper is to build the missing
bridge and to prove the theorems it implies.

\paragraph{The problem.}
We ask whether alignment dynamics can be modeled as variational free-energy minimization in
which a narrow misalignment intervention of strength $\lambda$ induces a phase transition in
the energy landscape, such that: (i) a critical threshold $\lambda^\ast$ separates robust
alignment from EM; (ii) the order of the transition is predictable; and (iii) the framework
predicts which intervention regimes yield a null result versus a sudden collapse, together with
quantitative precursors. The empirical literature poses exactly this and leaves it open,
remarking that the transition order ``may itself depend on the intervention regime''
\cite{nghiem2026trait,soligo2026em}.

\paragraph{Our approach and results.}
We model the misalignment order parameter $m\in\R$ (identified empirically with a toxic-persona
latent, the dominant principal component of trait drift, or a moral-susceptibility score) as
the internal coordinate of an FEP gradient flow $\dot m=-\partial_m F$. We then \emph{derive}
the effective potential from a two-component generative mixture: a Taylor expansion of the
FEP free energy about the aligned state yields the cusp-catastrophe normal form
\begin{equation}\label{eq:cusp-intro}
   F(m;a,h)=\tfrac14 b\,m^4+\tfrac12 a(\lambda)\,m^2-h(\lambda)\,m,\qquad b>0,
\end{equation}
with $a(\lambda)=a_0-c\lambda$ an alignment stiffness eroded by the intervention and
$h(\lambda)$ a directional bias set by its asymmetry. From the single family
\eqref{eq:cusp-intro} we prove:

\begin{itemize}[leftmargin=*,itemsep=2pt,topsep=2pt]
  \item \textbf{A critical threshold (Theorem~\ref{thm:threshold}).} The aligned equilibrium
    $m=0$ is asymptotically stable if and only if $a(\lambda)>0$, and loses stability when the
    Hessian eigenvalue crosses zero at the unique $\lambda^\ast=a_0/c$. Below $\lambda^\ast$:
    robust alignment. Above: emergent misalignment.
  \item \textbf{Order set by symmetry (Theorems~\ref{thm:pitchfork}--\ref{thm:fold}).}
    A symmetric intervention ($h\equiv0$) gives a supercritical pitchfork---a \emph{continuous}
    transition with mean-field exponent $\beta=\tfrac12$ (the gradual-LoRA regime). A biased
    intervention ($h\neq0$) unfolds the pitchfork into a fold (saddle-node), giving a
    \emph{catastrophic jump with hysteresis} on the exact cusp curve $27bh^2+4a^3=0$ (the
    full-finetune regime). This unifies the two empirical kinetic regimes in one normal form.
  \item \textbf{Early-warning laws (Theorem~\ref{thm:earlywarning}).} Near $\lambda^\ast$ the
    stationary variance and autocorrelation time of $m$ diverge (critical slowing down),
    giving a pre-failure precursor in which the internal order parameter \emph{leads} the
    behavioral jump.
\end{itemize}

\paragraph{Proof techniques.}
The threshold is a linear-stability calculation; the order of the transition follows from
center-manifold/Lyapunov--Schmidt reduction together with the gradient dynamic-transition
sign criterion of Ma--Wang \cite{mawang2008} and the catastrophic-bifurcation classification
of Kuehn \cite{kuehn2011}; the cusp curve is obtained by eliminating $m$ between $F'$ and
$F''$ (a resultant computation); the early-warning laws come from linearizing the governing
stochastic differential equation to an Ornstein--Uhlenbeck process. Every analytic claim is
checked symbolically (the resultant is recovered exactly as $b^2(4a^3+27bh^2)$) and numerically
(variance divergence, hysteresis folds matching the analytic curve to within $1\%$, and a
stationary law matching the Boltzmann form to a Jensen--Shannon divergence of
$6\times10^{-4}$ nats). Calibrated to the reported empirical thresholds, the model produces a
universality data-collapse and classifies null-versus-EM outcomes with ROC-AUC $0.996$ and
threshold MAE $0.056$.

\paragraph{Significance.}
The upshot is that \emph{when} an alignment intervention catastrophically fails is a computable
property of the landscape---the signs of $a(\lambda)$ and $h(\lambda)$---and the failure is
preceded by measurable critical-slowing-down signatures usable as an early-warning monitor.
This turns an empirical surprise into an a priori, monitorable risk calculation.

\paragraph{Organization.}
Section~\ref{sec:prelim} fixes notation, recalls the prerequisite results from the FEP and
critical-transition literature, and states the model. Section~\ref{sec:results} contains the
main theorems with complete proofs, illustrative examples, and the computational verification.
Section~\ref{sec:discussion} discusses implications, limitations, and open problems, and
Section~\ref{sec:conclusion} concludes.
